\chapter{What is established and what comes next}\label{ch:claims}

\section{The evidence ladder belongs inside the book}
We began with a clinic, a payment record and a physical question. We now have a declared state, a potential, an exact finite action account, a local cancellation proof, a group decomposition and an auditable way to compose transitions. The final task is to say what each of those achievements does and does not establish.

A project becomes misleading when a definition, a successful software check and a hope about society are reported in the same voice. The solution is not to weaken every sentence into uncertainty. It is to make the kind of statement clear. An identity can be exact under explicit assumptions. An empirical conclusion can be limited to its registered observations. An institutional proposal can be stated plainly as a proposal.

This edition is a theoretical and explanatory revision. Its arithmetic examples have been supplied for teaching and checked as such. They are not new scientific runs. Their purpose is to let a reader reconstruct the calculation and inspect the assumptions before any future result is interpreted.

\begin{center}\small
\begin{longtable}{P{0.23\linewidth}P{0.65\linewidth}}\toprule
Kind of statement & Example and appropriate interpretation\\\midrule
\endfirsthead
\toprule Kind of statement & Example and appropriate interpretation\\\midrule
\endhead
Definition & $V_i=(x_i-x_i^*)^2/(2\sigma_i^2)$ declares the model's ruler. It does not discover a law of nature.\\
Conditional identity & The exact quadratic expansion follows from that definition with fixed parameters.\\
Theorem under assumptions & Complete affected factors cancel; consistently linked endpoint differences telescope.\\
Implementation check & A named function matches a specified calculation at checked inputs. This requires its own code and test record.\\
Scientific observation & A registered execution produced particular outcomes under a frozen protocol.\\
Research hypothesis & A proposed actor/capacity mechanism may improve recovery under declared disturbances.\\
Institutional choice & An organization decides how to authorize action, support access or assign responsibility.\\\bottomrule
\end{longtable}
\end{center}

\section{The declared physical world}
The principal worked model contains one homogeneous, conserved stock distributed over finitely many locations. Stocks are nonnegative; an internal lossless transfer removes exactly what it adds elsewhere. Any further capacity or action constraint must be stated explicitly. Physical feasibility does not follow from a favourable potential value.

The reference vector is nonnegative and compatible with the fixed total stock. Every scale is positive and fixed during the account. These conditions make the Gaussian potential well defined and put the chosen reference in the basic fixed-mass domain. They do not prove that every action menu can reach it.

This small model deliberately omits many real features. It contains no automatic regeneration, no biological population threshold law and no restoring flow hidden inside the field definition. A future domain can include additional processes, but their physical carriers, state effects and assumptions have to be declared.

The simplification is useful because it exposes the accounting logic. It becomes misleading only if conclusions are transferred to the omitted features without justification. A clear theoretical model is an instrument for asking questions, not a miniature planet whose every property generalizes.

\section{The Gaussian potential and its exact consequences}
With fixed $x^*$ and positive $\sigma_i$,
\begin{align}
 V(x)&=\frac12\sum_i\left(\frac{x_i-x_i^*}{\sigma_i}\right)^2,\\
 \mu_i(x)&=\frac{x_i-x_i^*}{\sigma_i^2},\qquad
 H=\operatorname{diag}(1/\sigma_i^2).
\end{align}
The potential is nonnegative and zero exactly at the declared reference. Its gradient is a vector of local slopes. The Hessian specifies the curvature. All three objects are properties of the declared function.

For any fixed admissible increment $\delta$, expansion gives
\begin{equation}
 V(z+\delta)=V(z)+\mu(z)^T\delta+\frac12\delta^TH\delta.
\end{equation}
The expression is exact for the quadratic, not a short-step approximation. Positive definiteness of $H$ makes the quadratic term positive for nonzero $\delta$. Thus the initial tangent field reduction overstates the full field reduction by that term.

The general action equation is
\begin{equation}\label{eq:final-ebu}
 \boxed{E_a=V_{\rm loc}(z)-V_{\rm loc}(z+\delta_a)-C_a.}
\end{equation}
The principal examples set $C_a=0$. Where a residual process burden is used, it must describe declared physical burden, use compatible units and avoid duplicating represented effects. No price, wage or arbitrary penalty becomes physical burden merely by being placed in that term.

\section{The conditional laws in plain language}
First, a physical transfer must have a complete quantity account. In the main lossless model, the stock removed at one endpoint appears at another. A potential difference is not a substitute for this check.

Second, local valuation can be exact. Every unchanged factor has the same value before and after and therefore cancels. A coupled factor expands the required read neighbourhood; omitting it can even reverse the sign.

Third, a finite value includes curvature. A positive initial field does not guarantee a positive result for every finite quantity. The action must be evaluated at its complete endpoint, subject to the declared physical constraints.

Fourth, a group has a joint value. Fixed increments add in the state update, but their singleton values need not add through a nonlinear potential. The Gaussian cross-term formula explains the discrepancy without introducing extra physical stock or additional reward.

Fifth, common-path contributions sum to the group field difference under their stated assumptions. The sum follows from calculus. The interpretation of those contributions as physical path decompositions requires a model-realized path; their interpretation as ownership or capacity updates requires an additional rule.

Sixth, consecutive endpoint differences telescope under one fixed field. Closed undriven sequences have net value equal to minus the declared nonnegative process burden. External changes and revisions of the field definition require explicit additional terms.

\section{What the potential does not contain}
A potential tells us how a declared state is scored. It does not by itself choose an action, assign an actor, select an amount, specify a timescale or guarantee that a path is physically realizable. Each of those objects needs a model of its own.

In particular, the Gaussian shape is not evidence that actual states are normally distributed. The conserved total couples coordinates even if the Hessian is diagonal. A proposed probability law on the fixed-mass domain would need an explicit construction and empirical interpretation.

Nor does a reference automatically become a dynamic equilibrium. If a future transition rule leaves the reference unchanged, that establishes a particular invariance property. Attraction concerns what happens from other states. Homeostatic function concerns whether relevant systems continue operating under disturbances. These are related research questions, not alternative names for the minimum of a quadratic.

\section{The historical mathematics keeps its own scope}
The preserved Parts II and III contain earlier models with threshold penalties, flow laws, discrete-time bounds and preservation controllers. Their proofs and results remain statements about their specified definitions. This edition removes those mechanisms from the introductory core because they are not required to teach the general finite account.

It would be wrong to say that changing the introductory potential invalidates every historical theorem. It would be equally wrong to relabel old experimental findings as Gaussian findings. A theorem transfers when its assumptions hold; an observation transfers only through an additional justified argument about the relevant systems.

For example, unchanged-factor cancellation applies to the new separable quadratic by the same algebraic reasoning. A result that depended on an explicitly chosen threshold controller does not become a result for random-affordable actors merely because both programmes use the name EBU.

The book therefore retains the historical volumes as sources and explains the revised model directly. Compatibility notes and a chapter-disposition ledger accompany the manuscript. They identify where teaching responsibility moved or examples changed, rather than rewriting the past into a seamless but false narrative.

\section{A narrow historical evidence example}
The repository contains a finalized Gate 1D-C execution manifest recording thirty registered runs and six thousand trace rows. Its protocol, implementation identities and result artifacts define what that study tested. This is actual historical execution, distinct from the illustrations in this book.\cite{gate1dc}

The existence of those records matters: the project contains scientific work as well as prospective designs. Their existence does not establish that the new Gaussian mechanism has run, that all future harnesses have passed review, or that a real economy would behave as intended.

Detailed historical outcomes belong with their original comparisons and limitations. A short introductory paragraph cannot replace the frozen protocol or turn a positive finding on one outcome into general superiority. Readers who want to evaluate that evidence should follow the manifest and its identified sources.

\section{Why a software test is a different achievement}
A function that computes $4q-q^2/4$ correctly at chosen rational quantities can be checked without simulating a world. Such a check can catch a sign or coefficient error. It does not observe an actor selecting an action or a population recovering from disturbance.

A trajectory test is another class of work. It advances model states under declared rules. Its interpretation depends on the world, initialization, randomness, horizon, comparisons and outcomes fixed before the run. Calling it a unit test does not remove those scientific requirements.

Performance measurements are different again. Fast evaluation may make a large study practical, but speed is not evidence of model validity. A benchmark can establish a measured software property on a named host. It cannot establish whether the actions it would evaluate preserve homeostasis.

The distinction helps protect both kinds of work. Implementation checks can be celebrated for the errors they detect without being asked to support claims they never tested. Scientific evidence can be evaluated on its registered question without being confused with a build log.

\section{The current prospective question}
The new research direction asks what happens when actors choose by unweighted random selection among affordable, physically feasible actions or groups. Affordability is a separately declared rule based on persistent signed capacity accounts. The sampling unit and finite opportunity menu still need to be specified. EBU supplies the field calculation. It does not secretly choose the action that maximizes the value.

That separation makes the result scientifically open. A mechanism might recover, cycle, remain trapped, refuse useful actions or perform worse than a comparison rule. These are possible outcomes to distinguish, not defects to tune away after inspection.

The proposed study also needs a physical comparison arm using the same relevant opportunities without the capacity restriction. This would help identify the effect of the mechanism being studied. The comparator must be specified fairly; it should not be handicapped by a private information restriction or an artificially weaker physical menu.

No experimental parameter values, horizon or success threshold are chosen by this book. The examples teach equations. They do not quietly preregister a study because their numbers happen to be convenient.

\section{Decisions a study must make before execution}
A scientific protocol must define its event order. Does forcing occur before the action baseline? How are physically infeasible disturbances handled? How do the actions compose, and what is actually measured? Small changes can define different stochastic processes.

It must define the sampling unit and opportunity menu. Choosing an edge, an actor or a group uniformly need not produce the same distribution of physical actions. Duplicate descriptions of an identical action can change probabilities unless the menu convention addresses them. An empty affordable menu needs an explicit response.

The capacity proposal needs an owner map, a path interpretation, signed netting rules and a numerical policy for boundary comparisons. Group closure alone supplies none of these choices. They affect which actions can occur and therefore belong in the registered model.

Finally, the protocol needs outcome measures capable of distinguishing its claims. It should report physical validity, service where applicable, deviation, refusals and account behaviour without combining them into an unexplained winner label. A mechanism can improve one measure and worsen another.

\section{The same random seed is not always the same exposure}
Imagine two comparison arms receiving identical raw disturbance proposals. After their states diverge, one proposal may be feasible in one arm and infeasible in the other. If infeasible proposals are rejected, the applied physical disturbances differ even though the raw random numbers match.

This does not make paired comparison impossible. It means the coupling must be described precisely. The study can report both the shared proposals and the realized inputs, or use a design that makes the intended equality feasible. It should not claim identical physical exposure solely because a seed is shared.

The point is another example of the book's reading habit: identify the object. A random draw, an accepted disturbance and an observed state change are different objects. Scientific conclusions depend on which one was controlled.

\section{What would count as recovery?}
Recovery needs a declared reference condition, disturbance and time criterion. Does it mean exact return, approach within a specified neighbourhood, restored service, or a distribution concentrated near functional states? These are different propositions.

On a finite fixed-mass domain with positive fixed scales, deviation is bounded even under many unsuccessful mechanisms. Boundedness alone is therefore a weak discriminator. An idle or trapped system can remain bounded forever. The study needs a stronger outcome if its claim concerns useful restoration.

With continuing forcing, a physical state distribution might settle while cumulative audit or capacity accounts drift. Stationarity of the physical state is not stationarity of the augmented state containing all balances. Any claim about a long-run process must identify which coordinates it covers.

These distinctions should be made before results are observed. Otherwise it becomes too easy to describe whichever quantity looks stable as the intended homeostatic outcome.

\section{Calibration is a scientific and institutional task}
The field uses references and scales. Domain evidence must support their physical interpretation; institutional procedures must determine who may adopt or revise them. A mathematically convenient value is not automatically a defensible public standard.

Calibration can change whose actions receive positive or negative values. The reference determines the centre; the scale determines how strongly departures count relative to one another. A doubling of all scales divides this quadratic potential and its field values by four. That is a change in the numerical ruler, not a physical improvement.

Sensitivity analysis can show how conclusions depend on these choices. Independent measurements can constrain them. Appeals and review can address contested representation. None of these steps is replaced by high arithmetic precision.

\section{A revised clinic capstone}
Return to Vale with three linked needs: water, medicine storage and a pump repair. For the water demonstration, use $z=(18,2)\X$, references $(10,10)\X$, scales $(2,2)\X$ and an accepted lossless four-unit transfer. The endpoint is $(14,6)\X$. Conservation gives twenty units before and after, potential changes from sixteen to four, and field value is twelve.

Now add a separately reported delivery requirement of five units for the stated interval. The action supplies four, leaving one unmet on that requirement. The field improved and the service requirement was not fully met. Reporting both is more informative than assigning an unqualified success label.

The pump repair uses a material part and labour. Neither is silently converted into water stock. Its physical transformations need their own records. If the current water state contains no equipment-capacity coordinate, the present water potential does not issue an invented repair benefit. Later water movements can be observed under the repaired action map.

If medicine refrigeration also changes an energy account, preserve the separate resource coordinate. A water value of twelve does not automatically offset an energy burden of another numerical amount. The clinic's actual functioning depends on linked services, timing and quality, not solely on their scalar accounts.

Finally, a public institution might support the repair or guarantee access to medicine. That decision can be justified by need and evidence while leaving the physical record intact. The patient is neither the debit nor the balancing entry.

\section{Questions a reader should now be able to answer}
\textbf{Is EBU physical energy?} Not by the name alone. In this Gaussian model, the standardized potential is dimensionless. A physical application must declare its quantities and calibration. Energy is one possible represented carrier, not an automatic interpretation of every value.

\textbf{Is the scale a measurement error?} No. Here it is a declared width of the reference geometry. Sensor uncertainty is a separate statement about our knowledge of stock. It may affect a later uncertainty analysis but should not silently redefine the field.

\textbf{Does zero potential prove a functioning clinic?} It proves that the represented stocks match the declared references. A clinic also needs quality, timing, equipment, staff and access. Unrepresented requirements are not certified by that equality.

\textbf{Does positive marginal field prove positive finite value?} Only a sufficiently small movement may inherit the initial directional sign when the assumptions and process burden permit it. The finite equation includes curvature and the full endpoint.

\textbf{Must every action be sequential?} No. The book represents simultaneous groups. Their joint value is calculated at the combined endpoint. A historical theorem with a restricted execution rule retains that scope; it does not prohibit separate analytical treatment of groups.

\textbf{Does the path attribution settle ownership?} No. It provides a declared mathematical decomposition. Physical-path interpretation, causal identification, ownership and spendability need their own assumptions and authority.

\textbf{Can a negative action still be necessary?} Yes. The represented field and an urgent service requirement can point in different directions. The report should expose the conflict, and the relevant institution should decide how to respond.

\textbf{Can a positive action be refused?} Yes. It may be physically infeasible, unauthorized, unaffordable under a proposed capacity rule, or simply not selected. A value calculation is not a command.

\textbf{Why preserve a zero or negative example?} It tests whether we understand the equation rather than treating it as a reward label. Such examples can reveal missing state, overshoot, incomplete service or a genuine tradeoff.

\textbf{Can a field be improved later?} Yes. Scientific models should respond to evidence. Versioning and explicit model-change records prevent an improved description from being misreported as a physical action in the old history.

\textbf{What would justify a claim about an economy?} Evidence about measurement, behaviour, access, incentives, institutions and physical outcomes at the relevant scale. A finite accounting identity is one component, not the complete argument.

\section{A reader's final reconstruction exercise}
Choose a finite action with one or two physical edges. Write its source-side quantity, complete increment and predecessor state. Check feasibility and conservation before evaluating the potential. Declare references and positive scales with units and a compatible total.

Calculate the before and after potentials. Obtain the field value by subtraction and independently by the exact quadratic expansion. If there is a group, calculate the joint increment and compare with the same-baseline singleton sum. If you display path attributions, state what path they use and verify that they close.

Write a short receipt containing the equation version, physical inputs, calculation and evidence status. Add a separately stated service question. Then identify one human or institutional question your field does not answer. A complete response should allow another reader to reproduce the mathematics without guessing what the action or the claim was.

If a step remains uncertain, revisit the corresponding practice studio before drawing a scientific conclusion.

\section{The equation and the work ahead}
The core account can now be read without treating its symbols as a black box:
\[
 \boxed{\begin{aligned}\text{signed EBU}={}&\text{before potential}-\text{after potential}\\
 &-\text{declared residual burden}.\end{aligned}}
\]
The state and physical map say what changes. The potential says how that represented change is measured. Feasibility says which changes the physical model permits. An actor rule says what is chosen. Evidence says what happened. Institutions decide how people may use the account.

\begin{lawbox}{The foundation in one connected statement}
With a declared physical boundary, fixed field and complete action support, exact finite differences provide consistent local action values. Their group and sequential identities follow under explicit assumptions. Calibration, realized dynamics, useful recovery, fair access and an action-accounted economy remain substantive work to be demonstrated.
\end{lawbox}

Part II continues at Chapter 33: from the action's domain and derivations to
physical permission, software records, tests and historical evidence. The
question stays the same: what changed, under which assumptions, and what can
we conclude? A limitation, counterexample or failed comparison can be a useful
next result. Scientific integrity preserves those possibilities as carefully
as the equation.
